En una fàbrica, el 4\,\% de les bombetes surten defectuoses. En cada control s'agafen 20 bombetes a l'atzar
i es compten les defectuoses ($X$).

\begin{apartats}

\apartat[1,25]{0,75}
Justifica que $X$ segueix una distribució binomial i indica'n els paràmetres.

\begin{solucio}
Cada bombeta és defectuosa o no ho és (dos resultats), en cada control n'hi ha un nombre fix, 20, i la
producció és tan gran que la probabilitat de ser defectuosa és sempre $0{,}04$, independentment de les
altres bombetes. Per tant, $X\sim B(20;\,0{,}04)$.
\end{solucio}

\begin{tria}{binomial-tasca}
\itemtria{original}{1}{1,25}
Calcula la mitjana i la variància de $X$, i la probabilitat que en un control no hi hagi cap bombeta
defectuosa.

\begin{solucio}
$\mu=20\cdot0{,}04=0{,}8$ i $\sigma^2=20\cdot0{,}04\cdot0{,}96=0{,}768$.\\
$P(X=0)=0{,}96^{20}\approx0{,}442$.
\end{solucio}

\itemtria{parametres-inversos}{1}{1,25}
En una altra fàbrica, el nombre de bombetes defectuoses de cada control segueix una distribució binomial
de mitjana $1{,}2$ i variància $1{,}14$. Quantes bombetes s'agafen en cada control, i quin percentatge de
bombetes surten defectuoses?

\begin{solucio}
$np=1{,}2$ i $np(1-p)=1{,}14$. Dividint: $1-p=\dfrac{1{,}14}{1{,}2}=0{,}95$, i per tant $p=0{,}05$: surten
defectuoses el 5\,\% de les bombetes.\\
$n=\dfrac{1{,}2}{0{,}05}=24$ bombetes en cada control.
\end{solucio}
\end{tria}

\begin{nomesllarg}
\apartat{0,75}
Quin és el nombre de bombetes defectuoses més probable en un control? Justifica-ho.

\begin{solucio}
$P(X=0)\approx0{,}442$ i $P(X=1)=20\cdot0{,}04\cdot0{,}96^{19}\approx0{,}368$, i les següents són més
petites. El més probable és \textbf{cap}, tot i que la mitjana és $0{,}8$.
\end{solucio}
\end{nomesllarg}

\end{apartats}
